\section{Phase propagation and Gaussian controls}\label{S:nogo-phase}

A wave equation transports displacement and velocity together. Its scalar
solution value is not a state variable for composition. Likewise, a
Gaussian change of law transports an entire random solution, not merely
the solution evaluated at the mean of the new noise. These distinctions
explain the failures in this section. The scalar examples are diagnostics
for generic proof principles; their scope is stated separately from the
homogeneous stochastic wave model.

\subsection{NG-14: scalar moments do not compose across blocks}
\label{S:nogo-14}

Independent static time blocks look favorable for multiplying second
moments. Independence does remove correlations between their potentials,
but it does not reset the velocity to zero at a block boundary. The
original record computed formal coefficients only. We include the
integrable-series argument needed to turn that diagnostic into a
small-variance counterexample.

\setcounter{nogo}{13}
\begin{nogo}[Two independent wave blocks do not multiply scalar moments]
  \label{P:nogo-14}
  Let $Z_1,Z_2$ be independent centered Gaussians of variance $A$. Solve
  $y''=Z_1y$ on $[0,1]$ and $y''=Z_2y$ on $[1,2]$, with continuous phase
  at time one and initial phase $(y(0),y'(0))=(1,0)$. Then, as
  $A\downarrow0$,
  \[
    \E y(1)^2=1+\frac A3+O(A^2),\qquad
    \E y(2)^2=1+3A+O(A^2).
  \]
  In particular, $\E y(2)^2>(\E y(1)^2)^2$ for all sufficiently small
  positive $A$.
\end{nogo}
\begin{proof}
  Define the entire functions
  \[
    c(z)=\sum_{n\ge0}\frac{z^n}{(2n)!},\qquad
    s(z)=\sum_{n\ge0}\frac{z^n}{(2n+1)!}.
  \]
  The unit-time phase propagator and the terminal displacement are
  \[
    M(z)=\begin{pmatrix}c(z)&s(z)\\zs(z)&c(z)\end{pmatrix},\qquad
    y(2)=c(Z_2)c(Z_1)+s(Z_2)Z_1s(Z_1).
  \]
  Expanding through total degree two in $(z,w)$ gives
  \[
    c(w)c(z)+s(w)zs(z)
    =1+\frac32z+\frac12w
    +\frac5{24}z^2+\frac5{12}zw+\frac1{24}w^2
    +\text{terms of degree at least three}.
  \]
  After squaring, the coefficients of $z^2$ and $w^2$ are respectively
  $8/3$ and $1/3$. Independence and centering eliminate the $zw$ term.
  Similarly, the coefficient of $z^2$ in $c(z)^2$ is $1/3$.

  For completeness, these expansions admit the stated remainders after
  Gaussian expectation. The series of $c$ and $s$ have nonnegative
  coefficients, and
  \[
    |c(z)|+|s(z)|\le2e^{\sqrt{|z|}}.
  \]
  Consequently, the absolute series for the square of the displayed
  two-block displacement is bounded by
  $C(1+|z|)^2e^{2\sqrt{|z|}+2\sqrt{|w|}}$.
  This is integrable at any fixed Gaussian variance $A_1>0$. Absolute
  integration allows termwise expectation. Odd powers in either variable
  vanish; every remaining term has even total degree and a nonnegative
  coefficient. Thus, the expectation is a convergent series
  $\sum_{j\ge0}b_jA^j$ with $b_j\ge0$ at $A_1$. For $0\le A\le A_1$,
  \[
    \sum_{j\ge2}b_jA^j
    \le (A/A_1)^2\sum_{j\ge2}b_jA_1^j=O(A^2).
  \]
  The same argument applies to the one-block moment. Its square equals
  $1+2A/3+O(A^2)$, whereas the two-block moment equals
  $1+3A+O(A^2)$. Their difference is $7A/3+O(A^2)$.
  This proves No-go~\ref{P:nogo-14}.
\end{proof}

\begin{figure}
  \centering
  \begin{tikzpicture}[x=1cm,y=1cm]
    \node[ham/box,text width=1.8cm] (initial) at (0,0)
    {Initial phase\\$(1,0)$};
    \node[ham/box,text width=3cm] (middle) at (4.1,0)
    {After block 1\\$(c(Z_1),Z_1s(Z_1))$};
    \node[ham/box,text width=2.6cm] (final) at (8.7,0)
    {After block 2\\$M(Z_2)M(Z_1)e_1$};
    \draw[ham/arrow] (initial)--node[above,font=\small] {$M(Z_1)$} (middle);
    \draw[ham/arrow] (middle)--node[above,font=\small] {$M(Z_2)$} (final);
    \node[ham/note,text width=10cm] at (4.4,-1.7)
    {Retained displacement: $c(Z_2)c(Z_1)+s(Z_2)Z_1s(Z_1)$.\\
    Resetting the velocity would retain only $c(Z_2)c(Z_1)$.};
    \node[ham/note,text width=10cm] at (4.4,-3)
    {True second moment: $1+3A+O(A^2)$.\\
    Product of reset-block moments: $1+\frac23A+O(A^2)$.};
  \end{tikzpicture}
  \caption[Velocity survives a block boundary]{Two independent scalar static
    Gaussian blocks, each of variance $A$. The second block acts on both
    components of the first block's phase. This is a scalar diagnostic
    outside the target homogeneous-noise model; see
    No-go~\ref{P:nogo-14}.}
  \label{Fig:nogo-phase-blocks}
\end{figure}


This scalar model has spatially constant noise and independent static
blocks, outside Hypothesis~\ref{H:homog}. It excludes the unrestricted
scalar multiplication rule, not a phase-space operator estimate. Corollary~\ref{C:static-interaction} gives a valid replacement for independent
successive blocks: it controls both phase coordinates and includes explicit
block losses whose dyadic sum is finite. Its stated scope does not identify
products across overlapping levels with the original evolution.

\subsection{NG-15: the missing second-order mean term}\label{S:nogo-15}

Keeping the first random Duhamel term gives a nonnegative squared norm.
It is not the complete coefficient of the variance in the expected phase
energy. The mean of the second Duhamel term contributes at the same order.

\setcounter{nogo}{14}
\begin{nogo}[The first variance coefficient need not be nonnegative]
  \label{P:nogo-15}
  For the scalar phase propagator $M$ in No-go~\ref{P:nogo-14} and a
  standard Gaussian $Z$, write
  \[
    \E\|M(\sigma Z)x\|^2=\|M(0)x\|^2+\sigma^2x^{\mathsf T}Kx
    +O(\sigma^4).
  \]
  Then
  \[
    K=\begin{pmatrix}4/3&4/5\\4/5&17/45\end{pmatrix},\qquad
    (1,-2)K(1,-2)^{\mathsf T}=-\frac{16}{45}.
  \]
  The matrix $K$ is the full first coefficient in the variance parameter,
  not a covariance matrix. Omitting the second-order mean term also
  underestimates this coefficient for $x=(1,0)^{\mathsf T}$.
\end{nogo}
\begin{proof}
  The coefficients of $M(z)=M_0+zM_1+z^2M_2+\cdots$ are
  \[
    M_0=\begin{pmatrix}1&1\\0&1\end{pmatrix},\quad
    M_1=\begin{pmatrix}1/2&1/6\\1&1/2\end{pmatrix},\quad
    M_2=\begin{pmatrix}1/24&1/120\\1/6&1/24\end{pmatrix}.
  \]
  Gaussian expectation of the squared phase norm gives
  \[
    K=M_1^{\mathsf T}M_1+M_0^{\mathsf T}M_2+M_2^{\mathsf T}M_0.
  \]
  Matrix multiplication gives the stated entries and negative value.
  For $e_1=(1,0)^{\mathsf T}$, the first random contribution is
  $\|M_1e_1\|^2=5/4$, whereas $e_1^{\mathsf T}Ke_1=4/3$.
  Thus, it is not valid even as an upper bound for the full coefficient.

  The entire series are absolutely Gaussian integrable, including their
  squares, by the bound in No-go~\ref{P:nogo-14} with the additional
  polynomial factor from $zs(z)$. Absolute summation of the even-degree
  terms at a fixed nonzero $\sigma_1$ bounds the tail of degree at least
  four by $C\sigma^4$ for $|\sigma|\le|\sigma_1|$. Therefore, for
  $x=(1,-2)^{\mathsf T}$ the actual expected energy satisfies
  \[
    \E\|M(\sigma Z)x\|^2=5-\frac{16}{45}\sigma^2+O(\sigma^4).
  \]
  This proves No-go~\ref{P:nogo-15}.
\end{proof}

The example refutes arbitrary-phase covariance monotonicity for this
scalar flow. It says nothing against monotonicity of the original
solution-value moment with initial data $(1,0)$ in the homogeneous model.
Section~\ref{S:static-phase} retains both contributions to the coefficient.
The all-order estimate of Proposition~\ref{P:static-phase-all-orders} and
the interaction bound of Corollary~\ref{C:static-interaction} retain these
terms. Positivity of the first random variance does not control the cross
term with the mean.

\subsection{NG-16: replacing a tilted mean by a deterministic value}
\label{S:nogo-16}

The Gaussian density cost is exact. A loss may instead enter when the
random solution under the tilted law is replaced by a deterministic
surrogate. The following two calculations identify the loss without
asserting an obstruction to Gaussian shifts themselves.

\setcounter{nogo}{15}
\begin{nogo}[A deterministic surrogate can lose the annealed factor]
  \label{P:nogo-16}
  Let $A>0$ and $0<r<2$, and define
  \[
    J_c=\sup_{H\ge0}\{cAH^r-H^2/2\}.
  \]
  Then $J_c=c^{2/(2-r)}J_1$, whereas
  \[
    \sup_{H\ge0}\{2AH^r-H^2\}=2J_1<J_2.
  \]
  Moreover, for $Y=e^{\sigma Z}$, $Z$ standard Gaussian and $\sigma>0$,
  the full tilted-mean Cauchy--Schwarz bound recovers $\E Y^2$ exactly,
  while substituting $Y(h)=e^{\sigma h}$ loses a factor in its logarithm.
\end{nogo}
\begin{proof}
  Substitute $H=c^{1/(2-r)}K$ in the definition of $J_c$. Both terms
  acquire the factor $c^{2/(2-r)}$. The supremum $J_1$ is finite, since
  $r<2$, and positive, since for small positive $H$ the positive term
  dominates. Factoring two from the surrogate score proves the strict
  comparison because $2/(2-r)>1$.

  Let $Q_h$ be the Gaussian shift with density
  $L_h=e^{hZ-h^2/2}$ relative to the standard Gaussian law. Then
  $\E L_h^2=e^{h^2}$, so Cauchy--Schwarz gives
  \[
    \E Y^2\ge(\E_{Q_h}Y)^2e^{-h^2}.
  \]
  The exact tilted mean is $\E_{Q_h}Y=e^{\sigma^2/2+\sigma h}$.
  Maximizing its resulting exponent
  $\sigma^2+2\sigma h-h^2$ at $h=\sigma$ gives
  $e^{2\sigma^2}=\E Y^2$. Substitution of the deterministic value
  $Y(h)$ instead gives the maximum $e^{\sigma^2}$.
  This proves No-go~\ref{P:nogo-16}.
\end{proof}

These are generic scalar score calculations, not an SPDE no-go. They
identify which term is lost and show that retaining the full tilted mean
can avoid the loss. A deterministic surrogate remains useful if one
proves that its error is negligible on the intended exponential scale;
without such an estimate it cannot identify the annealed constant.

\subsection{NG-20: covariance profiles need not separate}
\label{S:nogo-20}

A fixed spatial profile times a scalar temporal factor is easier to
propagate than a general spacetime control. Such a product family is
legitimate as a restriction, but it does not represent every positive
control even in the white spatial model.

\setcounter{nogo}{19}
\begin{nogo}[A positive control need not have proportional covariance profiles]
  \label{P:nogo-20}
  Fix $0<a<1$ and one-dimensional spatial covariance
  $\Gamma=c_\Gamma\delta_0$, $c_\Gamma>0$. There is a nonnegative smooth
  compactly supported spacetime test $\psi$ for which
  \[
    h_\sigma(x)=c_\Gamma\int_0^1|\sigma-u|^{-a}\psi(u,x)\ud u,
    \qquad 0\le\sigma\le1,
  \]
  cannot be represented as $b(\sigma)U(x)$ with one fixed spatial $U$.
\end{nogo}
\begin{proof}
  Choose a nonzero nonnegative smooth $f_1$ supported in $(1/8,1/4)$,
  let $f_2(u)=f_1(1-u)$, and choose nonzero nonnegative smooth bumps
  $q_1,q_2$ supported respectively in $(-2,-1)$ and $(1,2)$. Set
  $\psi(u,x)=f_1(u)q_1(x)+f_2(u)q_2(x)$. Put
  \[
    A=\int_0^1u^{-a}f_1(u)\ud u,\qquad
    B=\int_0^1(1-u)^{-a}f_1(u)\ud u.
  \]
  Since $u<1-u$ on the support of $f_1$, we have $A>B>0$.
  Reflection in time gives
  \[
    h_0=c_\Gamma(Aq_1+Bq_2),\qquad
    h_1=c_\Gamma(Bq_1+Aq_2).
  \]
  Choose $x_i$ with $q_i(x_i)>0$. Disjoint supports imply
  \[
    h_0(x_1)h_1(x_2)-h_1(x_1)h_0(x_2)
    =c_\Gamma^2(A^2-B^2)q_1(x_1)q_2(x_2)>0.
  \]
  A representation by one $U$ would make this determinant zero.
  The white covariance has $\alpha=1$, so $1+a<2$ verifies the model
  parameter condition. This proves No-go~\ref{P:nogo-20}.
\end{proof}

This is an exact representation obstruction. It does not prove that
separable controls fail to attain a variational supremum; a theorem
reducing optimizers to a smaller family would be a different statement.
A useful weaker fact survives: because $|\sigma-u|^{-a}\ge1$ on the
unit time square, every nonnegative test has
$h_\sigma\ge\Gamma*\int_0^1\psi(u,\cdot)\ud u$.
This common lower profile does not by itself prove a phase-propagation
comparison. Reopening an exact product reduction requires extra structure
on the control, while an optimizer reduction requires a variational proof.

\subsection{NG-31: initial phase and the backward cone}\label{S:nogo-31}

Two seemingly small generalizations remove essential hypotheses: replacing
the specified initial phase by an arbitrary one, and checking agreement
of spatial cutoffs only on the terminal observation set. Simple
deterministic waves already distinguish these statements.

\setcounter{nogo}{30}
\begin{nogo}[Growth does not give every-phase lower bounds or local cutoff consistency]
  \label{P:nogo-31}
  The growth rate of a wave phase propagator need not be a lower growth
  rate for every nonzero initial phase. Also, two smooth initial-data
  masks that agree on an observation interval can give different free-wave
  displacements throughout that interval at a later time, even when their
  initial displacements agree everywhere.
\end{nogo}
\begin{proof}
  For the first assertion, let
  $B=\left(\begin{smallmatrix}0&1\\1&0\end{smallmatrix}\right)$.
  Its eigenvectors $(1,1)^{\mathsf T}$ and $(1,-1)^{\mathsf T}$ have
  eigenvalues $1$ and $-1$. Hence, $\|e^{tB}\|=e^t$ for $t\ge0$,
  but $e^{tB}(1,-1)^{\mathsf T}=e^{-t}(1,-1)^{\mathsf T}$.

  For the second assertion take the one-dimensional free wave with zero
  initial displacement and nonnegative smooth velocity $v$ supported in
  $(5,6)$, normalized by $\int v=1$. Choose smooth masks $\chi_s,\chi_l$
  between zero and one, with $\chi_s=1$ on $[-2,2]$ and support in
  $[-4,4]$, and $\chi_l=1$ on $[-14,14]$. Both masks agree on the
  observation interval $[-2,2]$. The masked velocities are nevertheless
  $\chi_sv=0$ and $\chi_lv=v$. D'Alembert's formula at time nine gives
  \[
    u_s(9,x)=0,\qquad
    u_l(9,x)=\frac12\int_{x-9}^{x+9}v(y)\ud y=\frac12,
    \qquad -2\le x\le2.
  \]
  Indeed, every interval $[x-9,x+9]$ in question contains $(5,6)$.
  This proves No-go~\ref{P:nogo-31}.
\end{proof}

The first example does not refute a terminal lower bound for the original
phase $(1,0)$. The second requires agreement of both displacement and
velocity on the full backward cone, not only on the observer. Finite
propagation with that guard survives. Thus, a successor must retain the
specified phase, or prove a suitable nonvanishing unstable component, and
must impose cutoff agreement on the relevant cone for both coordinates.
